\begin{frame}
\begin{itemize}
    \item On considère un fluide visqueux au repos que l'on suppose incompressible, 
homogène et ne subissant aucune force extérieure. Le problème de Stokes consiste à étudier 
la vitesse du fluide dans la situation suivante
\end{itemize}
\begin{figure}[!h]
    \centering
    \tdplotsetmaincoords{70}{120}
    \begin{tikzpicture}[tdplot_main_coords, scale=1.4]
        % ----- PLAQUE (plan y=0, dans le plan xz) -----
        \fill[blue!15, opacity=0.7] (-2.5,-2.5,0) -- (2.5,-2.5,0) -- (2.5,2.5,0) -- (-2.5,2.5,0) -- cycle;
        \draw[gray!60] (-2.5,-2.5,0) -- (2.5,-2.5,0) -- (2.5,2.5,0) -- (-2.5,2.5,0) -- cycle;
        % Label plaque
        \node[blue!60!black] at (2.8, 1.5, 0) {plaque $y=0$};

        % ----- REPÈRE -----
        % Axe x (e1)
        \draw[->, thick] (0,0,0) -- (2.5,0,0) node[anchor=north] {$x$};
        \draw[->, red] (0,0,0) -- (1,0,0) node[above left] {$e_1$};
        % Axe y (e2) vers le haut
        \draw[->, thick] (0,0,0) -- (0,0,2.5) node[anchor=south] {$y$};
        \draw[->, green] (0,0,0) -- (0,0,1) node[right] {$e_2$};
        % Axe z (e3)
        \draw[->, thick] (0,0,0) -- (0,2.5,0) node[anchor=west] {$z$};
        \draw[->, blue] (0,0,0) -- (0,1,0) node[below right] {$e_3$};
        % Origine
        \filldraw[black] (0,0,0) circle (1.5pt) node[below] {$O$};

        % ----- FLUIDE (demi-espace y>0) -----
        \fill[cyan!10, opacity=0.3]
            (-2.5,-2.5,0) -- (2.5,-2.5,0) -- (2.5,-2.5,2.2)
            -- (-2.5,-2.5,2.2) -- cycle;
        \fill[cyan!10, opacity=0.3]
            (2.5,-2.5,0) -- (2.5,2.5,0) -- (2.5,2.5,2.2)
            -- (2.5,-2.5,2.2) -- cycle;
        \fill[cyan!10, opacity=0.3]
            (-2.5,-2.5,2.2) -- (2.5,-2.5,2.2) -- (2.5,2.5,2.2)
            -- (-2.5,2.5,2.2) -- cycle;
        \fill[cyan!10, opacity=0.3]
            (2.5,2.5,0) -- (-2.5,2.5,0) -- (-2.5,2.5,2.2)
            -- (2.5,2.5,2.2) -- cycle;
        \node[cyan!50!black] at (2.2, -1.5, 1.8) {fluide $y>0$};

        % ----- FLÈCHE DE DÉPLACEMENT DE LA PLAQUE -----
        \draw[->, line width=2pt, red!80!black]
            (-0.5, -1.8, 0) -- (1.2, -1.8, 0)
            node[midway, above, sloped] {\small $V e_1$};

        % ----- POINT X -----
        \filldraw[black] (-0.4, 1.0, 1.2) circle (1.5pt);
        \node[above right] at (-0.4, 1.0, 1.2) {$X = (x,y,z)$};
    \end{tikzpicture}
    \caption{Configuration du premier problème de Stokes}
    \label{fig:pbmStokes}
\end{figure}
\end{frame}

\begin{frame}
\begin{itemize}
\item On considère donc la vitesse du fluide
\begin{eqnarray*}
    \bfm{u} : \bbR^3 \times \bbR_+ \longrightarrow \bbR^3, \
    (X,t) \longmapsto \big(u_1(X,t), u_2(X,t), u_3(X,t)\big)
\end{eqnarray*} 
\item Le contexte du premier problème de Stokes nous amène à quelques 
simplifications. 
\begin{equation*}
    \bfm{u} (X,t) = \big(u_1(y,t),0,0\big)
\end{equation*} 
\item Comme le fluide est incompressible et homogène on a 
\begin{eqnarray*}
    \rho(X,t) \equiv \rho > 0 \ \ \ &\text{ et }& \ \ \ \Div(\bfm{u})=0
\end{eqnarray*} 
\item Les équations du mouvement du milieu continu nous donne 
\begin{equation}\label{PFD1}
    \rho \frac{D\bfm{u}}{Dt} = \Div(\sigma) + \rho \bfm{f}
\end{equation}
\end{itemize}
\end{frame}

\begin{frame}
\begin{itemize}
    \item Pour un fluide de second grade, le tenseur des contraintes est 
donné par le théorème de Rivlin-Ericksen 
\begin{equation*}
    \sigma := -pI + \mu A_1 + \alpha_1 A_2 + \alpha_2 A_1^2
\end{equation*}
où
\begin{itemize}
    \item[$\triangleright$] $pI$ est la contrainte isotropique avec $p$ la pression,
    \item[$\triangleright$] $A_1 := \nabla \bfm{u} + \nabla^t \bfm{u}$,
    \item[$\triangleright$] $A_2 := \dfrac{D}{Dt}A_1 + A_1\nabla \bfm{u} + \nabla^t\bfm{u} A_1$,
    \item[$\triangleright$] $\mu >0$ est la viscosité du fluide,
    \item[$\triangleright$] $\alpha_1$ et $\alpha_2$ sont des paramètres thermodynamiques qui doivent vérifier 
    \begin{eqnarray*}
        \alpha_1 \geq 0 \ \ \ \ \ \text{ et } \ \ \ \ \ \alpha_1 + \alpha_2 = 0
    \end{eqnarray*}
    pour que le modèle soit physiquement admissible.
\end{itemize}
\end{itemize}
\end{frame}

\begin{frame}
\begin{itemize}
    \item Le tenseur des contraintes est
\begin{equation*}
    \sigma = \begin{pmatrix}
        -p+\alpha_2(\partial_yu_1)^2 & \mu\partial_yu_1 + \alpha_1 \partial_{ty}u_1 & 0 \\
        \mu \partial_yu_1+\alpha_1\partial_{ty} u_1 & -p+(2\alpha_1+\alpha_2)(\partial_yu_1)^2 & 0 \\
        0 & 0 & -p
    \end{pmatrix}
\end{equation*} 
\item On applique l'opérateur divergence à $\sigma$ et on injecte dans \eqref{PFD1} et 
\begin{equation*}
    \left\{\begin{array}{rcl}
        \rho\partial_tu_1 & = & -\partial_xp + \mu \partial_{yy}u_1 + \alpha_1 \partial_{yty}u_1  \\
        0 & = & - \partial_yp+2(2\alpha_1+\alpha_2)\partial_yu_1\partial_{yy}u_1 \\
        0 & = & - \partial_zp
    \end{array}\right.
\end{equation*}
    \item L'invariance par translation dans les directions $e_1$ et $e_3$ implique que 
la pression ne dépend pas de $x$ et $z$, donc $\partial_xp=0$.
\item On obtient finalement l'équation de Rayleigh-Stokes homogène en une dimension
\begin{equation*}
    \partial_t u_1 - \nu \partial_{yy}u_1 - \alpha \partial_{tyy}u_1=0
\end{equation*}
où $\nu := \dfrac{\mu}{\rho}$ et $\alpha := \dfrac{\alpha_1}{\rho}$.
\end{itemize}
\end{frame}

%\begin{frame}
%\begin{itemize}
%    \item L'invariance par translation dans les directions $e_1$ et $e_3$ implique que 
%la pression ne dépend pas de $x$ et $z$, donc $\partial_xp=0$.
%\item On obtient finalement l'équation de Rayleigh-Stokes homogène en une dimension
%\begin{equation*}
%    \partial_t u_1 - \nu \partial_{yy}u_1 - \alpha \partial_{tyy}u_1=0
%\end{equation*}
%où $\nu := \dfrac{\mu}{\rho}$ et $\alpha := \dfrac{\alpha_1}{\rho}$.
%\end{itemize}
%\end{frame}